% This is samplepaper.tex, a sample chapter demonstrating the
% LLNCS macro package for Springer Computer Science proceedings;
% Version 2.21 of 2022/01/12
%
\documentclass[runningheads]{llncs}

\usepackage[T1]{fontenc}
% T1 fonts will be used to generate the final print and online PDFs,
% so please use T1 fonts in your manuscript whenever possible.
% Other font encondings may result in incorrect characters.
%

\usepackage{graphicx}
\usepackage{epsfig}
\usepackage{subcaption}
\usepackage{calc}
\usepackage{amssymb}
\usepackage{amstext}
\usepackage{amsmath}
\usepackage{multicol}
\usepackage{pslatex}
\usepackage{algorithm2e}
\usepackage[bottom]{footmisc}
\usepackage{booktabs}
\usepackage{multirow}
\usepackage{svg}

% Used for displaying a sample figure. If possible, figure files should
% be included in EPS format.
%
% If you use the hyperref package, please uncomment the following two lines
% to display URLs in blue roman font according to Springer's eBook style:
%\usepackage{color}
%\renewcommand\UrlFont{\color{blue}\rmfamily}
%\urlstyle{rm}
%
\begin{document}
%
\title{PAMR: Link Prediction on Knowledge Graphs via Path-Aware Multi-modal Reasoning}
%
\titlerunning{Link Prediction on KG via Path-Aware Multi-modal Reasoning}
% If the paper title is too long for the running head, you can set
% an abbreviated paper title here
%
\author{Hai-Dang Huynh-Lam\inst{1,2}\orcidID{0000-0001-9487-7281} \and
  Ngoc-Thao Nguyen\inst{1,2}\orcidID{1111-2222-3333-4444} \and
Ngoc-Thanh Le\inst{1,2}\orcidID{0009-0009-6426-2652}}
%
\authorrunning{Hai-Dang et al.}
% First names are abbreviated in the running head.
% If there are more than two authors, 'et al.' is used.
%
\institute{
  Faculty of Information Technology, University of Science, Ho Chi Minh City, Vietnam
  % \email{lncs@springer.com}\\
  % \url{http://www.springer.com/gp/computer-science/lncs}
  \and
  Vietnam National University, Ho Chi Minh City, Vietnam\\
\email{\{hlhdang,nnthao,lnthanh\}@fit.hcmus.edu.vn}}
%
\maketitle  % typeset the header of the contribution
%
\begin{abstract}
  The abstract should briefly summarize the contents of the paper in
  150--250 words.

  \keywords{Knowledge Graph \and Multi-modal Learning \and Path-based Reasoning \and Link Prediction \and Transformers.}
\end{abstract}
%
%
%

\section{\uppercase{Introduction}}
\label{sec:introduction}

% TODO: Expand this section with full introduction narrative.
% Should include: problem motivation, background on KGs and link prediction,
% challenges of multimodal reasoning, contributions list.

\paragraph{PAMR.}
The proposed framework, \textbf{Path-Aware Multimodal Reasoning (PAMR)}, addresses
link prediction on multimodal knowledge graphs by scoring structural paths between
query entities using a Transformer encoder over frozen KGE token embeddings.

\paragraph{Contributions.}
PAMR introduces two key design choices: (i)~a paired sampling scheme that, unlike
conventional uniform negative sampling, yields batches with exactly balanced positive
and negative examples per query, and (ii)~a token representation that concatenates
real and imaginary relation embeddings while using only the fused real entity
representation. These design choices enable the Transformer to learn discriminative
path representations from frozen KGE token embeddings without modifying the underlying
knowledge graph representations.

\section{\uppercase{Related Work}}
\label{sec:related-work}

The development of the PAMR framework lies at the intersection of several core research areas in artificial intelligence, including graph-based relational reasoning, multimodal representation learning, and sequential modeling. In this section, we survey the literature across three foundational dimensions to establish the methodological context and justify the necessity of PAMR.

\subsection{Path-Based Reasoning on Knowledge Graphs}
\label{sec:rw-path-reasoning}

Path-based reasoning on knowledge graphs extends conventional single-hop link prediction by extracting and evaluating multi-hop relational sequences that connect source and target entities. Foundational methods in this domain relied heavily on statistical random walks, such as the Path Ranking Algorithm (PRA)~\cite{lao2011random}. The transition to deep learning introduced architectures like NeuralLP~\cite{yang2017differentiable} and RNNLogic~\cite{qu2021rnnlogic}, which formulate path extraction as differentiable logic programming or probabilistic modeling, respectively, enabling end-to-end learning of logical rules. Concurrently, deep reinforcement learning (RL) was adapted to navigate graph spaces. DeepPath~\cite{xiong2017deeppath} and MINERVA~\cite{das2018go} model pathfinding as a Markov Decision Process, where an agent learns to step toward the target entity using shaped rewards. To mitigate the sparse reward issues common in long-path traversal, DIVA~\cite{chen2018variational} framed the problem as variational inference, while CURL~\cite{zhang2022learning} proposed a multi-agent hierarchical RL framework.

More recently, graph neural networks (GNNs) and sequence-to-sequence models have been employed to efficiently encode path structures. RED-GNN~\cite{zhang2022knowledge} and NBFNet~\cite{zhu2021neural} parameterize the Bellman-Ford algorithm with neural components to capture generalized path representations through message passing. AdaProp~\cite{zhang2023adaprop} improves upon this by adaptively filtering propagation paths, while SQUIRE~\cite{bai2022squire} introduces a Transformer-based sequence-to-sequence framework that translates queries directly into paths. However, most path-based methods rely on learned rather than frozen embeddings, use reinforcement learning or message-passing rather than Transformers, and do not handle multimodal entities. This interwining between embedding and the sequential modeling process increases computational burden and limiting the ability to leverage advanced representations.

\subsection{Multimodal Knowledge Graph Embeddings}
\label{sec:rw-mmkge}

Multimodal knowledge graph embedding (MMKGE) seeks to overcome the structural sparsity of traditional KGs by injecting auxiliary modalities, such as images and text, into the representation learning process. Early approaches, such as IKRL~\cite{xie2017image} and TransAE~\cite{wang2019transae}, demonstrated that aligning visual and textual features with structural embeddings could resolve entity ambiguity and improve link prediction. The MMKG dataset collection~\cite{liu2019mmkg} further established baselines for numeric and visual entity matching. Subsequent models focused on filtering modality noise and handling structural complexity; for instance, RSME~\cite{wang2021rsme} uses a forget-gate mechanism to discard conflicting multimodal signals, while OTE~\cite{tang2020orthogonal} maps multimodal entities via orthogonal relation transforms to preserve topological distances. Additionally, MMRNS~\cite{xu2022multimodal} incorporates cross-modal attention to enhance negative sampling with multimodal semantics.

State-of-the-art MMKGE architectures increasingly rely on Transformers. MKGformer~\cite{chen2022mkgformer} employs a hybrid multi-level fusion approach to integrate visual and textual representations. MoSE~\cite{zhao2022mose} introduces a modality-split framework to prevent interference during fusion, and VISTA~\cite{lee2023vista} extends multimodal encoding directly to the triplet level using specialized decoding Transformers. Finally, VL-KGE~\cite{efthymiou2026vlkge} leverages vision-language models to handle modality asymmetry in large-scale graphs. Notably, existing multimodal KGE methods fuse modalities into a single entity embedding but do not expose multimodal token sequences to a path-aware Transformer for reasoning, as the fusion primarily serves triple scoring rather than path discrimination. Therefore, no widely used benchmark treats multimodal information as entities on the graph that can connect to different conceptual entities rather than being bound to a single entity. This multimodal-as-entity setting is what the proposed method evaluates on.

\subsection{Transformers for Knowledge Graph Reasoning}
\label{sec:rw-transformers}

The application of Transformers and pre-trained language models to KGs leverages self-attention mechanisms to contextualize relational data. KG-BERT~\cite{yao2019kg} pioneered this by treating triples as natural language sequences for sequence classification. To address the computational bottleneck of cross-encoding every candidate, StAR~\cite{wang2021star} partitions triples into asymmetric context representations. Moving beyond purely textual encoding, CoKE~\cite{wang2019coke} tokenizes graph structures, passing edge and path sequences through a Transformer encoder to learn context-adaptive embeddings. HittER~\cite{chen2021hitter} extracts features through a hierarchical architecture, aggregating local neighborhood information using masked entity prediction.

For complex logical queries, KG-Transformer~\cite{liu2022mask} utilizes a masked pre-training strategy coupled with a Mixture-of-Experts sparse activation. Generative and prompt-based paradigms have also emerged, with KGT5~\cite{saxena2022kgt5} formulating link prediction as a sequence-to-sequence task, GenKGC~\cite{zhu2022genkgc} utilizing relation-guided prompts for target generation, and CSPrompter~\cite{chen2023csprompter} bridging the gap between symbolic logic and large language models via consistent structure prompting. However, Transformer-based KG methods typically encode triples or subgraphs rather than interleaved path sequences, use language-model pretraining rather than frozen KGE embeddings, and mostly target static text KGs rather than multimodal graphs, let alone the multimodal graphs under multimodal-as-entity setting.

% \subsection{Evaluation Approaches}
% \label{sec:rw-evaluation}

% Evaluation of link prediction models spans both training signal quality
% and test protocol rigor. Margin-based ranking losses, widely used in
% this domain, rely heavily on negative sampling quality. Traditional
% uniform random sampling~\cite{bordes2013translating} or Bernoulli
% distributions~\cite{wang2014knowledge} often yield easily
% distinguishable or false-negative samples, while self-adversarial
% weighting~\cite{sun2019rotate} and generative adversarial frameworks
% like KBGAN~\cite{cai2018kbgan} improve gradient informativeness at
% the cost of training stability. More efficient alternatives such as
% NSCaching~\cite{zhang2019nscaching} and CAKE~\cite{niu2022cake}
% achieve adversarial-level performance through cached hard negatives
% and commonsense priors, respectively. However, these methods are
% designed for triple-level discrimination and do not guarantee balanced
% positive-to-negative ratios per batch, making them unsuitable for
% path-level paired sampling.

% Standardized benchmarks complement these training-side advances by
% providing rigorous test protocols. WN18RR~\cite{dettmers2018convolutional}
% and FB15k-237~\cite{toutanova2015observed} eliminate test-set leakage
% and inverse relation bias, while
% YAGO3-10~\cite{mahdisoltani2013yago3} evaluates scalability across
% multilingual data. With the shift towards multimodality,
% domain-specific graphs like WikiArt-MKG~\cite{efthymiou2026vlkge}
% offer heterogeneous structures spanning artworks, artists, and textual
% attributes. Recent surveys~\cite{niu2025survey,liang2024survey} have
% cataloged these developments by task orientation and graph type.
% Nevertheless, standard benchmarks are either purely structural without
% multimodal features, or lack explicit path annotations for
% supervision, leaving WikiArt-MKG as the only dataset that combines
% both properties yet has not been used for Transformer-based path
% reasoning.

Taken together, existing work leaves several open questions that PAMR addresses: how to explicitly execute path-aware reasoning on heterogeneous multimodal entities, how to parse and project frozen complex-valued KGEs into a Transformer's token space without initialization deadlock, and how to stabilize sequence ranking via a perfectly balanced query-level paired sampling mechanism.

\section{Method}
\label{sec:method}
\subsection{Overview}
\label{sec:overview}

A KG is a directed multi-relational graph
$\mathcal{G} = (\mathcal{V}, \mathcal{E}, \mathcal{R})$ where
$\mathcal{V}$ is a set of entities, $\mathcal{R}$ a set of relation types,
and $\mathcal{E} \subseteq \mathcal{V} \times \mathcal{R} \times \mathcal{V}$
a set of labelled edges. Each edge $(s, r, o) \in \mathcal{E}$
denotes that subject entity $s$ is connected to object entity $o$
via relation $r$.

Path-based reasoning focuses on using structural paths and connection
patterns within KGs to infer new relations or answer queries. Given
a query $(s, r, ?)$, the target is to approximate the conditional
distribution $P(\hat{o} \mid s, r, \mathcal{G})$ over candidate
objects $\hat{\mathbb{O}}$. By using the shortest paths $p_o$ between $s$ and $\hat{o}$ in $\mathcal{G}$ as the most important factor contributing to the distribution,
the target reduces to $P(p_o \mid s, r)$. A scoring function
$\hat{f}_\theta$ is learned to predict this target and then assigns a real-valued score to each path given the
query relation $r$, afterward the candidate with the highest-scoring path is selected as the answer.

In a MMKG, entities are associated with heterogeneous modalities, for example artworks are inherently visual, while artists,
styles, genres, and locations are characterized textually. We model
this as a multimodal-as-attachment setting, where each entity $v \in \mathcal{V}$
possesses modality-specific embeddings $\mathbf{v}_i$ indexed by
$i \in \mathbb{M}$, with $\mathbb{M}$ denoting the modality set.
The unified embedding $\mathbf{v}$ is obtained
by aggregating over available modalities: $\mathbf{v} =
\text{AGGREGATE}_{i \in \mathbb{M}} (\mathbf{v}_i)$. In this proposal, we formulate the MMKG setting of modality-as-entity where each modality information is converted to a separate entity, this creates a mapping $\mathcal{M}: \mathcal{V} \to \mathbb{M}$ that maps each entity to its original modality.

The proposed method, PAMR, consists of four components. First, a frozen
pre-trained KGE model supplies token embeddings for entities and relations.
Second, a bidirectional BFS extracts shortest paths between query pairs under the modality-as-entity setting.
Third, a Transformer encoder scores each path by contextualising
interleaved entity and relation tokens. Fourth, a margin-based ranking
loss trains the scorer on paired positive and negative paths. Each
component is detailed in the following subsections.

\subsection{Embedding Module}
\label{sec:embedding-module}

To supply pre-encoded relational structure without incurring additional
training cost, we use a frozen pre-trained KGE model.  A KGE model
learns entity embeddings $\mathbf{E} \in \mathbb{H}^{|\mathcal{V}|}$
and relation embeddings $\mathbf{R} \in \mathbb{H}^{|\mathcal{R}|}$
such that a scoring function $h: \mathcal{V} \times \mathcal{R} \times
\mathcal{V} \to \mathbb{R}$ assigns higher values to true triples
than to false ones.  The embedding space $\mathbb{H}$ is a
$D$-dimensional continuous manifold.

In a multimodal KG, entities possess heterogeneous modalities.
VL-KGE~\cite{efthymiou2026vlkge} constructs a unified embedding for
every entity by fusing its available modalities into a single point
in $\mathbb{H}$.  For each entity $v \in \mathcal{V}$, let
$\mathcal{M}(v) \subseteq \mathbb{M}$ denote the set of available
modalities.  The unified entity embedding $\mathbf{v}$ is obtained
by averaging over available modalities:
$\mathbf{v} = \frac{1}{|\mathcal{M}(v)|} \sum_{i \in \mathcal{M}(v)} \mathbf{v}_i$,
where $\mathbf{v}_i$ is the modality-specific embedding for modality
$i \in \mathbb{M}$.  This per-entity averaging naturally handles
modality asymmetry, where different entity types carry different
modality combinations without requiring a fixed global fusion scheme.

The effective dimensionality of entity and relation representations
depends on the KGE architecture: complex-valued models output $2D$
dimensions while real-valued models output $D$. The Transformer
hidden dimension is set to $d_{\text{model}} = \max(d_{\mathcal{V}},
d_{\mathcal{R}})$, and two learnable linear projections align entity
and relation token dimensions to $d_{\text{model}}$. These
projections are learned during PAMR training while the KGE model
remains frozen.

\subsection{Path Extraction}
\label{sec:path-extraction}

To provide the Transformer with explicit, interpretable connectivity evidence rather than
relying solely on the implicit proximity captured by KGE embeddings, we extract shortest
paths between query pairs via bidirectional breadth-first search (BFS), bounded by a
maximum hop count $K$.

\paragraph{Multimodal and non-multimodal configurations.}
We define two graph configurations that produce qualitatively different path spaces.
The multimodal configuration $\mathcal{G}_{\text{multimodal}}$ comprises the complete
graph with all entity types, where paths may bridge across different entity types,
exploiting cross-type relationships for reasoning.  The concept-only configuration
$\mathcal{G}_{\text{concept}}$ restricts the graph to conceptual entities and
edges where both endpoints are conceptual, confining paths to a homogeneous
subgraph that excludes multimodal information.  The latter corresponds to
Path-Aware Reasoning (PAR), which relies solely on structural connectivity
without multimodal signals.

Formally, a path $p$ between $s$ and $\hat{o}$ comprises $\ell + 1$ edges, where
$\ell$ denotes the hop count bounded by $K$.
Paths with $\ell \leq 0$ are excluded from training and evaluation as they
carry no structural reasoning signal, which include unreachable pairs and
trivial self loops.

\subsection{PAMR Architecture}
\label{sec:pamr-architecture}

\subsubsection{Token Interleaving}
\label{sec:token-interleaving}

% TODO: remove overused text bold in verbal text, use math notations that are used in prior sections, embedding notations should re-use the embedding module before rather than let new definition here. replace \mathbf{r}_q with \mathbf{r} only as this only has single query relation r, use d instead of d_model as we use same dimension as prior embedding, L rather than L_tr. Remove the mentioning of node type embeddings also. The length of position encoding is based on max hop count K defined in earlier subsection where the length of vectorized path associated with this K hops is 2K + 3.
To encode path structure as a flat sequence, entity and relation tokens are interleaved
with a classification token and the query relation prepended.
Let $\mathbf{e}^{\text{raw}}_i \in \mathbb{R}^{d_{\mathcal{V}}}$ be the raw KGE embedding
of entity $i$ and $\mathbf{r}^{\text{raw}}_j \in \mathbb{R}^{d_{\mathcal{R}}}$ the raw
relation embedding of relation $j$.
Projected token embeddings are:
\begin{equation}
  \mathbf{v}_i = \operatorname{Proj}_{\mathcal{V}}(\mathbf{e}^{\text{raw}}_i),\qquad
  \mathbf{r}_j = \operatorname{Proj}_{\mathcal{R}}(\mathbf{r}^{\text{raw}}_j)
  \label{eq:token-proj}
\end{equation}
both lying in $\mathbb{R}^{d_{\text{model}}}$.
A path with $\ell$ hops comprises $\ell+1$ entity tokens and $\ell$ edge tokens.
These are interleaved into a token sequence of length $L = 2\ell + 1$:
\begin{equation}
  \mathbf{T} = \bigl[\mathbf{c};\; \mathbf{r}_q;\;
    \mathbf{v}_{\mathbf{s}};\; \mathbf{r}_{1};\; \mathbf{v}_{2};\; \mathbf{r}_{2};\;
  \mathbf{v}_{3};\; \ldots;\; \mathbf{v}_{\mathbf{o}}\bigr]
  \in \mathbb{R}^{L \times d_{\text{model}}}
  \label{eq:token-seq}
\end{equation}
Here $\mathbf{c} \in \mathbb{R}^{d_{\text{model}}}$ is a learnable classification token
at position 0, initialised from $\mathcal{N}(0, 0.02^2)$.
$\mathbf{r}_q$ is the query
relation representation at position 1, while entity tokens $\mathbf{v}_{\mathbf{s}}, \mathbf{v}_{2}, \ldots, \mathbf{v}_{\mathbf{o}}$
occupy positions $2 + 2i$ for $i = 0, \ldots, \ell$, and edge tokens $\mathbf{r}_{1}, \ldots, \mathbf{r}_{\ell}$
occupy positions $3 + 2i$ for $i = 0, \ldots, \ell - 1$.
Entity positions receive additive \textbf{node type embeddings}
$\boldsymbol{\tau}_i \in \mathbb{R}^{d_{\text{model}}}$ from a learned embedding table
with one entry per node type plus a padding zero.
All positions receive \textbf{absolute positional embeddings}
$\mathbf{p}_i = \operatorname{PE}(i)$ via a learned embedding
$\operatorname{PE}: \{0, \ldots, L-1\} \to \mathbb{R}^{d_{\text{model}}}$.
The final token input to the Transformer is:
\begin{equation}
  \mathbf{T}^{\text{in}}_i = \mathbf{T}_i \;+\; \mathbf{p}_i \;+\; \boldsymbol{\tau}_{\tau(i)} \cdot \mathbb{1}_{\mathcal{V}}(i)
  \label{eq:token-in}
\end{equation}
Here $\mathbb{1}_{\mathcal{V}}(i) = 1$ iff position $i$ is an entity position.
\subsubsection{Transformer Encoder}

\label{sec:transformer-encoder}
% TODO: use d' instead of d_ff
The token sequence is contextualised through a stack of $L_{\text{tr}}$
Transformer encoder layers~\cite{vaswani2017attention} with $h$ attention heads each:
\begin{equation}
  \mathbf{H}^{(0)} = \mathbf{T}^{\text{in}},\qquad
  \mathbf{H}^{(l)} = \operatorname{EncLayer}^{(l)}\bigl(\mathbf{H}^{(l-1)},\; \mathbf{M}\bigr)
  \quad (l = 1, \ldots, L_{\text{tr}})
  \label{eq:transformer}
\end{equation}
Here $\mathbf{M} \in \{0, 1\}^{L}$ is a padding mask ($\mathbf{M}_i = 1$ indicates an
invalid position) for paths shorter than the batch maximum $N_{\max}(\mathcal{B})$.
Each encoder layer applies:
\begin{align}
  \tilde{\mathbf{H}} &= \operatorname{LayerNorm}\bigl(
  \mathbf{H}^{(l-1)} + \operatorname{MHA}(\mathbf{H}^{(l-1)}; \mathbf{M})\bigr)
  \label{eq:attn}\\
  \mathbf{H}^{(l)} &= \operatorname{LayerNorm}\bigl(
  \tilde{\mathbf{H}} + \operatorname{FFN}(\tilde{\mathbf{H}})\bigr)
  \label{eq:ffn}
\end{align}
with $\operatorname{FFN}(\mathbf{x}) = \mathbf{W}^{(2)} \operatorname{GELU}(\mathbf{W}^{(1)} \mathbf{x} + \mathbf{b}^{(1)}) + \mathbf{b}^{(2)}$
and inner dimension $d_{\text{ff}} = 4 \times d_{\text{model}}$.
GELU~\cite{hendrycks2016gaussian} activation
and Post-LN layer normalisation are used throughout.

\subsubsection{Readout and Scoring}

\label{sec:readout}

% TODO: convert from paragraph-based explanation to more verbal one where words are used instead of \paragraph{} commands, merge the block equation below into single inline equation by making w_2^T GELU(...) + b_2.
To obtain a fixed-dimensional path representation, the contextualised token
sequence is aggregated via one of two readout strategies from
$\mathbf{H}^{(L_{\text{tr}})} \in \mathbb{R}^{L \times d_{\text{model}}}$:

\begin{description}
  \item[CLS.] $\mathbf{z} = \mathbf{H}^{(L_{\text{tr}})}[0, :] \in \mathbb{R}^{d_{\text{model}}}$, the classification token at position 0.
  \item[Mean.] $\mathbf{z} = \frac{\sum_{i=0}^{L-1} \mathbf{H}^{(L_{\text{tr}})}[i, :] \cdot (1 - \mathbf{M}[i])}{\sum_i (1 - \mathbf{M}[i])}$, averaging over valid tokens.
\end{description}

A two-layer MLP maps $\mathbf{z}$ to a scalar score $s \in \mathbb{R}$:
\begin{equation}
  \mathbf{h} = \operatorname{GELU}\bigl(\mathbf{W}_1 \mathbf{z} + \mathbf{b}_1\bigr),
  \qquad
  s = \mathbf{w}_2^{\top} \mathbf{h} + b_2
  \label{eq:readout}
\end{equation}
with dropout $p = 0.1$ applied after GELU.
All linear layers use Xavier
uniform initialisation and zero bias.

\subsection{Training}

\label{sec:training}

To train a path scorer with stable supervision, each query must contribute equal gradient
signal regardless of its number of negative samples.
We construct a paired sampling
dataset that guarantees both balance and fairness.
\subsubsection{Paired Sampling}

\label{sec:paired-sampling}

For each query $q$, the query generation process produces one
\textit{positive} path $p^+_q$, which connects $\mathbf{s}$ to the
true object $\mathbf{o}$, and $K$ \textit{negative} paths
$\mathcal{S}^-_q = \{p^-_{q,1}, \ldots, p^-_{q,K}\}$, which connect
$\mathbf{s}$ to sampled false objects. To construct the paired
dataset $\mathcal{D}_{\text{pair}}$, we iterate over each query $q$
and pair every negative path $p^- \in \mathcal{S}_q^-$ with its
corresponding positive $p^+_q$, appending each pair to an index
$\mathcal{I}$. The total number of pairs equals $\sum_q
|\mathcal{S}_q^-|$, ensuring that every negative is visited exactly
once per epoch. Collation interleaves positive and negative pairs
into flat sequences:
$\operatorname{batch}(B) = [p^+_0, p^-_0, p^+_1, p^-_1, \ldots, p^+_{B-1}, p^-_{B-1}]$.
Here $B$ is the query batch size. Each training batch contains
exactly $B$ positive and $B$ negative samples, producing a guaranteed
50/50 composition regardless of the underlying positive to negative
ratio. Variable length paths are zero-padded to the batch maximum
$N_{\max}(\mathcal{B}) = \max_{p \in \mathcal{B}} (\ell_p + 1)$
entity tokens, with padding masks $\mathbf{M}$ as defined in
Section~\ref{sec:token-interleaving}.
\subsubsection{Loss Function}
\label{sec:loss-functions}

PAMR is trained with a margin-based ranking loss that groups samples
by query identifier and enforces a margin between the positive score
and all negative scores for that query. Let $\mathcal{B} = \{(s_1, y_1), \ldots, (s_{2B}, y_{2B})\}$
denote a training batch with $y_i \in \{0, 1\}$. The loss is defined as:
\begin{equation}
  \mathcal{L}_{\text{rank}} = \frac{1}{|\mathcal{Q}_{\mathcal{B}}|} \sum_{q \in \mathcal{Q}_{\mathcal{B}}}\; \frac{1}{|S_q^+||S_q^-|} \sum_{s^+ \in S_q^+} \sum_{s^- \in S_q^-}
  \max\bigl(0,\, m - s^+ + s^-\bigr)
  \label{eq:loss-rank}
\end{equation}
Here $\mathcal{Q}_{\mathcal{B}}$ denotes the set of unique query identifiers in $\mathcal{B}$, and
$m = 1.0$ is the margin. This loss directly optimises the ranking objective and prevents
score collapse.

\section{Experiment}
\label{sec:experiment}
\subsection{Implementation Details}
\label{sec:implementation}

All experiments run on a Vast.ai instance equipped with an NVIDIA
GeForce RTX 5060 Ti and an AMD EPYC 7763 processor allocated with 8 threads
and 16\,GB virtual RAM. The software stack comprises Python~3.11,
PyTorch~2.7.1 with CUDA~12.8 and cuDNN~9.

\paragraph{Path extraction.}
The bidirectional BFS of Section~\ref{sec:path-extraction} is
implemented in C++ (GNU GCC 13.2, compiled with C++23 standards)
and invoked from Python. The graph adjacency list is loaded once
into memory, and queries are distributed across 8 threads for
parallel path search. Per-thread outputs are merged into a single
path table per split. Precomputed paths are cached on disk as
pickle files for fast reuse across training runs.

\paragraph{GPU memory management.}
To accommodate the large frozen embedding table alongside the trainable
Transformer, GPU memory is managed through expandable memory segments
and periodic cache defragmentation after model initialisation.

\paragraph{Paired sampling implementation.}
The paired dataset is constructed once per epoch from precomputed
path tables, with per-query pairing parallelised across CPU cores.
Variable length paths are zero-padded to the batch maximum and
collated with padding masks as described in
Section~\ref{sec:paired-sampling}.

\paragraph{Training configuration.}
PAMR is trained with AdamW~\cite{loshchilov2019decoupled} at learning
rate $\eta = 0.01$, weight decay $\lambda = 0.01$, $\beta_1 = 0.9$,
$\beta_2 = 0.999$. The ranking loss (Eq.~\ref{eq:loss-rank}) is
used with margin $m = 1.0$. Batch size is 8 queries (16 samples per
batch with the guaranteed 50/50 positive/negative composition).
Training proceeds for 50 epochs. The Transformer encoder uses
$L_{\text{tr}} = 8$ layers and $h = 8$ attention heads with
$d_{\text{ff}} = 2048$ and dropout $p = 0.1$. The hidden dimension
$d_{\text{model}} = 1536$ is auto-detected from the frozen VL-KGE
model (ComplEx, $D = 768$). Only PAMR parameters are updated; the
frozen VL-KGE model remains untouched.

\subsection{Datasets}
\label{sec:datasets}

Experiments are conducted on WikiArt-MKG~v2~\cite{efthymiou2026vlkge},
a large-scale multimodal knowledge graph of fine art containing
224,166 entities and 22 relation types. Artworks are represented
visually through image embeddings, while attributes such as styles,
genres, years, and locations are represented textually. This modality
asymmetry requires models to handle entities with different
combinations of available modalities. The dataset comprises 224,166
entities, 216,564 artworks and 7,602 attributes, connected by 22
relation types, totalling 7,877,220 training triples, 208,513
validation triples, and 208,368 test triples. The graph is split
chronologically: training data precedes validation data, which in
turn precedes test data. Negative samples are generated via rejection
sampling with 1 negative per training query and 10 negatives per
validation/test query. For evaluation, three relation types
introduced as preprocessing inverses
and 10 negatives per validation/test query. For evaluation, three
relation types introduced as preprocessing inverses
(\texttt{isRelatedToArtwork}, \texttt{isRelatedToArtist},
\texttt{hasCreatedArtwork}) are excluded, leaving the four
relations between artists (\texttt{isInfluencedBy},
\texttt{isInfluencedOn}, \texttt{isPupilOf}, \texttt{isTeacherOf})
as primary evaluation targets. These four relations are chosen
because they connect artist entities, enabling evaluation under both
multimodal and non-multimodal reasoning settings. In the complete
graph, paths may traverse artworks and attribute entities such as
styles, genres, years, and locations, exercising the model's ability
to align visual and textual modalities. In the artist-only subgraph,
artwork entities are excluded from path extraction while the same
candidate set is retained, providing a controlled comparison that
isolates the contribution of multimodal path structure.

Path extraction uses bidirectional BFS with max hops $H = 5$, base
fanout $F = 100$, and decay factor $\lambda = 0.5$
(Section~\ref{sec:path-extraction}). Two graph configurations are
evaluated: $\mathcal{G}_{\text{full}}$ (the complete graph with all
entity types) and $\mathcal{G}_{\text{artist}}$ (the artist-only
subgraph restricted to edges between artists).

\subsection{Evaluation Metrics}
\label{sec:evaluation-metrics}

For each query, the model assigns a scalar score to every candidate
object and ranks them in descending order. We report Mean Reciprocal
Rank (MRR) and Hit-rate at cutoff $K$ (Hits@$K$), computed both
globally and disaggregated by relation type.

\subsection{Results}
\label{sec:results}

% TODO: splits into paragraphs including Comparative Study and Per-Relation Analysis, compress each parts to half the length with selection on more essential analysis among the existing one.
Table~\ref{tab:results} reports macro-averaged MRR and Hits@$K$ across
the four evaluation relations between artists. All methods share the same
frozen ComplEx embeddings ($D = 768$, $d_{\text{model}} = 1536$) to
ensure a fair comparison. The VL-KGE baselines score triples directly
without path reasoning, while PAR uses only structural paths and PAMR
additionally exploits multimodal token representations.

\begin{table*}[t]
  \centering
  \setlength{\tabcolsep}{2pt}
  \renewcommand{\arraystretch}{1.15}
  {\footnotesize
    \begin{tabular}{l|cccc|cc}
      \toprule
      \multirow{2}{*}{Metric} &
      \multicolumn{4}{|c|}{VL-KGE Embedding Baselines} &
      \multicolumn{2}{|c|}{Ours} \\
      \cmidrule(lr){2-5} \cmidrule(lr){6-7}
      & TransE & RotatE & DistMult & ComplEx & PAR & PAMR \\
      \midrule
      MRR  & 0.3196 & 0.3272 & 0.3107 & \underline{0.3630} & 0.2212 & \textbf{0.3789} \\
      Hits@1 & 0.2832 & 0.1824 & 0.3181 & \textbf{0.3870} & 0.0576 & \underline{0.3455} \\
      Hits@3 & 0.4938 & 0.4188 & 0.5283 & \textbf{0.6245} & 0.1987 & \underline{0.5759} \\
      Hits@5 & 0.7032 & 0.5973 & 0.7046 & \underline{0.7719} & 0.3306 & \textbf{0.7819} \\
      \bottomrule
  \end{tabular}}
  \caption{Macro-averaged ranking metrics across four evaluation
    relations between artists. \textbf{Bold} indicates the best per metric;
  \underline{underlines} indicate the second best.}
  \label{tab:results}
\end{table*}

Table~\ref{tab:results} reveals a consistent ordering across method
families. The results naturally stratify into three tiers: the four
VL-KGE embedding baselines occupy a moderate performance band, the
path-aware reasoning model, PAR, falls below this band, and the full
multimodal path-aware model, PAMR, surpasses it. This ordering holds
across all four metrics, with the strongest baseline, ComplEx, and PAMR
trading first and second positions depending on the cutoff.

Among the embedding baselines, ComplEx consistently outperforms the
other three architectures, achieving 0.363~MRR and 0.625~Hits@3
compared to the 0.31--0.33~MRR range of TransE, RotatE, and DistMult.
This advantage is expected: ComplEx's Hermitian dot product captures
asymmetric relations more effectively than the translation-based
TransE, the diagonal DistMult, or the complex rotation of RotatE, all
of which struggle with the diverse relation types in the art domain
graph. TransE, RotatE, and DistMult cluster tightly together,
differing by less than 0.02~MRR, suggesting that beyond ComplEx, the
choice of scoring function has limited impact on this particular graph
structure.

The path-aware reasoning model (PAR) underperforms all four VL-KGE
baselines, reaching only 0.221~MRR and 0.058~Hits@1. This gap is
substantial: PAR scores nearly 40\% below the weakest baseline,
DistMult at 0.311~MRR, and 73\% below ComplEx at Hits@1. The result
indicates that explicit path topology alone, without access to
multimodal entity features, provides weaker signal than the implicit
relational proximity encoded in the KGE embeddings. In the artist-only
subgraph $\mathcal{G}_{\text{artist}}$, paths consist solely of
artist-to-artist edges, which lack the rich visual and textual
attributes present on artwork and attribute entities. Without these
multimodal tokens, the Transformer has limited discriminative capacity
to distinguish factual from counterfactual paths.

Adding multimodal token representations through PAMR closes this gap
and surpasses the strongest baseline. PAMR achieves the best MRR
at 0.379, a 4.4\% improvement over ComplEx, and the best Hits@5 at 0.782,
a 1.3\% improvement. At the stricter cutoffs, PAMR remains
competitive with ComplEx: Hits@1 reaches 0.345 versus ComplEx's 0.387,
and Hits@3 reaches 0.576 versus ComplEx's 0.625. The gains at MRR
and Hits@5, combined with competitive performance at Hits@1 and Hits@3,
suggest that PAMR's path-based scoring improves the average rank of the
true object across the candidate set, even when it does not always
place the correct answer at the very top.

The crossover between MRR/Hits@5, where PAMR leads, and Hits@1/Hits@3,
where ComplEx leads, reveals a qualitative difference in scoring
behavior. ComplEx's direct triple scoring concentrates probability
sharply on a small set of candidates, occasionally placing the true
object at rank~1 but with less consistency across the full ranking.
PAMR's path-based scoring distributes probability more broadly, which
improves the average reciprocal rank and deeper recall metrics while
occasionally ceding the top position. This pattern aligns with the
design goal of path-aware reasoning: the model leverages explicit
structural evidence to make robust ranking decisions across the entire
candidate set, rather than relying on a single high-confidence
prediction. The multimodal tokens—image embeddings for artworks and
text embeddings for attributes—provide the Transformer with richer
input signals than the structural embeddings alone, enabling
finer-grained discrimination between factual and counterfactual paths.

To complement the macro-averaged results, Figures~\ref{fig:mrr},
\ref{fig:hits1}, \ref{fig:hits3}, and~\ref{fig:hits5} present the
per-relation breakdown for each metric. The four relations exhibit
distinct performance profiles, indicating that the benefit of
multimodal path reasoning depends on the relational context.

\begin{figure*}[t]
  \centering
  \includesvg[width=\linewidth]{figures/MRR.svg}
  \caption{Per-relation MRR across all methods.}
  \label{fig:mrr}
\end{figure*}

\begin{figure*}[t]
  \centering
  \includesvg[width=\linewidth]{figures/Hits1.svg}
  \caption{Per-relation Hits@1 across all methods.}
  \label{fig:hits1}
\end{figure*}

\begin{figure*}[t]
  \centering
  \includesvg[width=\linewidth]{figures/Hits3.svg}
  \caption{Per-relation Hits@3 across all methods.}
  \label{fig:hits3}
\end{figure*}

\begin{figure*}[t]
  \centering
  \includesvg[width=\linewidth]{figures/Hits5.svg}
  \caption{Per-relation Hits@5 across all methods.}
  \label{fig:hits5}
\end{figure*}

Two relations strongly favour PAMR. On \texttt{isPupilOf}, PAMR
achieves 0.390~MRR and 0.909~Hits@5, the highest Hits@5 across all
methods and relations, surpassing ComplEx by 0.038~MRR and
0.091~Hits@5. On \texttt{isInfluencedOn}, PAMR attains 0.415 MRR,
exceeding ComplEx by 0.090~MRR, the largest single-relation gap in
the experiment. These relations connect artists through stylistic
lineage and creative influence, where shared visual motifs and textual
descriptions of artistic movements provide rich multimodal context.
The Transformer can exploit these signals when scoring paths, enabling
finer discrimination between factual and counterfactual continuations
than ComplEx's direct triple scoring.

In contrast, \texttt{isTeacherOf} and \texttt{isInfluencedBy} favour
ComplEx. On \texttt{isTeacherOf}, ComplEx achieves 0.422~MRR versus
PAMR's 0.390, and on \texttt{isInfluencedBy}, ComplEx reaches
0.354~MRR versus PAMR's 0.321. These relations encode more abstract
or indirect associations: teacher-of relationships may involve
institutional or pedagogical connections that are not strongly
reflected in visual or textual attributes, while influence-by
relationships can span distant artistic periods with limited shared
multimodal context. In such cases, ComplEx's direct semantic
proximity in the embedding space captures the relational signal more
efficiently than the structural paths available in the graph.

PAR remains the weakest method across all four relations, confirming
the earlier finding that path structure alone is insufficient. Its
performance is particularly poor on \texttt{isPupilOf} at 0.209~MRR and
0.030~Hits@1 and on \texttt{isInfluencedBy} at 0.185~MRR and 0.035~Hits@1,
where the artist-only subgraph provides minimal discriminative paths
without the multimodal attributes carried by artwork and attribute
entities. The variation in the PAMR-to-ComplEx gap across relations
aligns with the design goal of multimodal path reasoning: the approach
is most beneficial when the relation connects entities with rich shared
visual or textual context, and least beneficial when the relation
encodes abstract or infrequent associations where direct embedding
proximity suffices.

\begin{credits}
  \subsubsection{\ackname} A bold run-in heading in small font size at the end of the paper is
  used for general acknowledgments, for example: This study was funded
  by X (grant number Y).

  \subsubsection{\discintname}
  It is now necessary to declare any competing interests or to specifically
  state that the authors have no competing interests. Please place the
  statement with a bold run-in heading in small font size beneath the
  (optional) acknowledgments\footnote{If EquinOCS, our proceedings submission
  system, is used, then the disclaimer can be provided directly in the system.},
  for example: The authors have no competing interests to declare that are
  relevant to the content of this article. Or: Author A has received research
  grants from Company W. Author B has received a speaker honorarium from
  Company X and owns stock in Company Y. Author C is a member of committee Z.
\end{credits}
%
% ---- Bibliography ----
%
% BibTeX users should specify bibliography style 'splncs04'.
% References will then be sorted and formatted in the correct style.
%
\bibliographystyle{splncs04}
\bibliography{refs}

% \begin{thebibliography}{8}
% \bibitem{ref_article1}
% Author, F.: Article title. Journal \textbf{2}(5), 99--110 (2016)

% \bibitem{ref_lncs1}
% Author, F., Author, S.: Title of a proceedings paper. In: Editor,
% F., Editor, S. (eds.) CONFERENCE 2016, LNCS, vol. 9999, pp. 1--13.
% Springer, Heidelberg (2016). \doi{10.10007/1234567890}

% \bibitem{ref_book1}
% Author, F., Author, S., Author, T.: Book title. 2nd edn. Publisher,
% Location (1999)

% \bibitem{ref_proc1}
% Author, A.-B.: Contribution title. In: 9th International Proceedings
% on Proceedings, pp. 1--2. Publisher, Location (2010)

% \bibitem{ref_url1}
% LNCS Homepage, \url{http://www.springer.com/lncs}, last accessed 2023/10/25
% \end{thebibliography}
\end{document}
